\documentclass[11pt,a4paper]{article}

\usepackage[T1]{fontenc}
\usepackage[utf8]{inputenc}
\usepackage{mathptmx}
\usepackage{amsmath,amssymb,amsthm}
\usepackage{amsfonts}
\AtBeginDocument{\let\hbar\relax
  \DeclareMathSymbol{\hbar}{\mathord}{AMSb}{"7E}}
\usepackage{geometry}
\usepackage{microtype}
\usepackage{hyperref}
\usepackage{booktabs}
\usepackage{array}
\usepackage{enumitem}
\usepackage{setspace}
\usepackage{titlesec}
\usepackage{abstract}
\usepackage{natbib}

\geometry{top=1.1in, bottom=1.1in, left=1.15in, right=1.15in}

\titleformat{\section}{\normalfont\bfseries\large}{\thesection}{1em}{}
\titleformat{\subsection}{\normalfont\bfseries\normalsize}{\thesubsection}{1em}{}
\titleformat{\subsubsection}{\normalfont\bfseries\itshape\normalsize}{\thesubsubsection}{1em}{}
\titlespacing*{\section}{0pt}{14pt}{4pt}
\titlespacing*{\subsection}{0pt}{10pt}{3pt}
\titlespacing*{\subsubsection}{0pt}{8pt}{2pt}

\renewcommand{\abstractnamefont}{\normalfont\bfseries}
\renewcommand{\abstracttextfont}{\normalfont\small}
\setlength{\absleftindent}{0.5cm}
\setlength{\absrightindent}{0.5cm}

\theoremstyle{plain}
\newtheorem{theorem}{Theorem}[section]
\newtheorem{proposition}[theorem]{Proposition}
\newtheorem{corollary}[theorem]{Corollary}
\theoremstyle{definition}
\newtheorem{definition}[theorem]{Definition}
\theoremstyle{remark}
\newtheorem{remark}[theorem]{Remark}

% Shortcuts
\newcommand{\Pact}{\mathcal{P}_{\mathrm{act}}}
\newcommand{\ARA}{\mathcal{A}_{\mathrm{RA}}}
\newcommand{\ket}[1]{|#1\rangle}
\newcommand{\bra}[1]{\langle #1|}
\newcommand{\braket}[2]{\langle #1 | #2 \rangle}

\begin{document}

\title{\textbf{Actualization as Physical Mechanism: \\
Implications of Relational Actualism for \\
Quantum Computing and Quantum Thermodynamics}}

\author{Joshua F.\ Sandeman \\
\small Independent Researcher \\
\small \texttt{sansuikyo@gmail.com}}

\date{March 2026}

\maketitle
\begin{center}
  {\small\textit{Preprint:}~\href{https://doi.org/10.5281/zenodo.19198285}
    {doi.org/10.5281/zenodo.19198285}\par}
\end{center}

\begin{abstract}
Relational Actualism (RA) proposes that the transition from quantum potentia to
classical actuality is grounded in a precise physical criterion: an irreversible
increase in quantum relative entropy $\Delta S(\rho\|\sigma_0) > 0$ with respect
to the vacuum. In the perturbative regime this coincides with the on-shell
dispersion relation $p^\mu p_\mu = m^2c^2$, at which a mediating boson crosses
the kinematic threshold from virtual to real and an irreversible vertex is
inscribed into the growing causal directed acyclic graph (DAG) of the universe. This paper develops the implications of this framework---established in
the companion papers RAQM, RAGC, RACI, and RAHC~\citep{Sandeman2026RAHC}---for two fields that have lacked a
precise physical account of irreversibility at the quantum level: quantum computing
and quantum thermodynamics. For quantum computing, the Kinematic Snap provides a
principled physical floor on decoherence, distinguishes genuinely off-shell
(decoherence-immune) interaction regimes from on-shell ones, and gives the
threshold theorem a physical rather than purely operational interpretation. For
quantum thermodynamics, RA grounds Landauer's principle in the causal DAG,
resolves the Maxwell's demon paradox without information-theoretic postulates,
gives the fluctuation theorems a precise actualization-event interpretation, and
establishes that the thermodynamic arrow of time is structural rather than
emergent. A key technical result: the quantum relative entropy---the central
quantity in both quantum information theory and RA's actualization
criterion---is frame-independent by a compiled Lean~4 theorem
(\texttt{frame\_independence} in \texttt{RA\_AQFT\_Proofs\_v10.lean}), proved
using the continuous functional calculus unitary-conjugation lemma from the
Lean-QuantumInfo library~\citep{Meiburg2025LQI}. A key quantitative connection: the Kinematic Coherence Bound
threshold $p_{\mathrm{th}} \approx 10^{-2}$ is shown to be consistent with the
galactic-scale coupling constant $\xi$ of the companion paper
RADM~\citep{Sandeman2026RADM} to within ${\sim}2\%$ using the full baryonic disk mass~\citep{BlandHawthorn2016}, at the Planck scale, via the
thermodynamic field equations established in RAGC~\citep{Sandeman2026RAGC}.
The Kinematic Coherence Bound threshold $p_{\mathrm{th}} \approx 10^{-2}$
also enters the Poisson-CSG generative measure of RAGC
(equation for $c_k^{\mathrm{RA}}$): it is the per-link actualization
probability that determines whether the causal DAG is locally 3D or 2D,
connecting quantum error correction thresholds to galactic-scale
dimensional reduction and the Tully-Fisher relation~\citep{Sandeman2026RADM}.
The fault-tolerance threshold, the speed of light $c = l_P/t_P$, and
the biological Causal Firewall are all the same Erd\H{o}s-R\'{e}nyi
percolation transition at different scales of the causal graph. The
Einstein field equations are shown to be the unique consistent
macroscopic description of the RA-CSG by Lovelock's theorem, with GR
exact at $\mu \gg 1$ and RA's departures from GR (rotation curves,
Hubble tension, Planck corrections) arising precisely where $\mu \sim 1$.
Throughout, implications are
stated conditionally on the correctness of the RA framework, which is currently
under peer review.

\medskip\noindent
\textbf{Keywords:} Relational Actualism \textperiodcentered{} Kinematic Snap
\textperiodcentered{} Quantum Computing \textperiodcentered{} Decoherence
\textperiodcentered{} Quantum Thermodynamics \textperiodcentered{} Landauer's
Principle \textperiodcentered{} Maxwell's Demon \textperiodcentered{} Fluctuation
Theorems \textperiodcentered{} Quantum Relative Entropy
\textperiodcentered{} Formal Verification
\end{abstract}

% ─────────────────────────────────────────────────────────────────────────────
\section{Introduction}
\label{sec:intro}

Two of the most practically consequential fields in contemporary physics share a
common foundational gap: neither quantum computing nor quantum thermodynamics
possesses a precise physical account of when, and by what mechanism, a quantum
process becomes irreversible. Quantum computing treats decoherence as an
engineering challenge to be suppressed, characterised by phenomenological
timescales $T_1$ and $T_2$ whose physical origins are described only
statistically. Quantum thermodynamics invokes entropy, work, and irreversibility
while relying on the same unitary, time-reversible quantum mechanics that makes
those concepts difficult to ground.

Relational Actualism (RA)~\citep{Sandeman2026a,Sandeman2026RAGC,Sandeman2026RACI}
proposes that this gap can be closed by a single physical criterion: actualization
occurs when a mediating boson crosses the on-shell dispersion relation
$p^\mu p_\mu = m^2c^2$. Below this threshold, exchange is virtual, off-shell,
and reversible. At and above it, the exchange is real, on-shell, and permanently
inscribed into the causal DAG. This is the \emph{Kinematic Snap}---the precise
physical mechanism by which quantum potentia becomes classical actuality.

The present paper develops the implications of this framework for quantum computing
(Section~\ref{sec:computing}) and quantum thermodynamics
(Section~\ref{sec:thermodynamics}), and then discusses the technical connection
to quantum information theory via the quantum relative entropy formalism
(Section~\ref{sec:qit}). All implications are stated conditionally: they follow
\emph{if} the RA actualization criterion is correct. The empirical predictions
in Section~\ref{sec:predictions} are the appropriate test.

We note at the outset that implications of this kind are not merely
interpretational. The Kinematic Snap criterion differs from decoherence-based
accounts in a precise, experimentally distinguishable way: it identifies a sharp
threshold (Prediction P3 of RAQM~\citep{Sandeman2026a}) rather than a continuous
crossover. That distinction has practical consequences for both fields developed
here.

% ─────────────────────────────────────────────────────────────────────────────
\section{Implications for Quantum Computing}
\label{sec:computing}

\subsection{The Physical Floor on Decoherence}

Standard quantum error correction characterises decoherence by the phenomenological
timescales $T_1$ (energy relaxation) and $T_2$ (phase coherence), which are
treated as empirical parameters to be minimised through engineering. RA gives these
parameters a physical interpretation and identifies the mechanism that sets them.

In RA, decoherence occurs when the system exchanges on-shell bosons with its
environment. Specifically, the actualization timescale is
\begin{equation}
  \tau_{\mathrm{act}} \;\sim\; \frac{\hbar}{\Delta E_{\mathrm{env}}},
  \label{eq:tact}
\end{equation}
where $\Delta E_{\mathrm{env}}$ is the characteristic energy available for
on-shell boson emission into the environmental degrees of freedom
(RAQM Prediction~P2~\citep{Sandeman2026a}). This is not a statistical parameter;
it is the time required for an emitted on-shell boson to propagate outward and
entangle irreversibly with the environmental boundary.

The practical implication is immediate: \emph{the engineering target for
decoherence suppression is $\Delta E_{\mathrm{env}}$, not the coupling strength
alone}. Reducing the energy available in the environment for on-shell exchange
directly extends $\tau_{\mathrm{act}}$. This gives a physical interpretation to
the well-known empirical strategy of operating superconducting qubits at millikelvin
temperatures: reducing $k_B T$ reduces $\Delta E_{\mathrm{env}}$, suppressing
on-shell exchange with thermal phonons. RA identifies this as the mechanism rather
than a phenomenological correlation.

Furthermore, \eqref{eq:tact} implies a physical floor that cannot be overcome by
any engineering improvement that does not reduce $\Delta E_{\mathrm{env}}$. No
amount of improved substrate quality, shielding, or control pulse optimisation
will suppress decoherence below $\hbar / \Delta E_{\mathrm{env}}$ if the
environmental energy scale is not itself reduced. This is a stronger statement
than the standard decoherence-theoretic account, which predicts continuous
improvement in principle.

\subsection{The Virtual/Real Distinction as a Decoherence-Immune Regime}

A non-trivial implication of the Kinematic Snap concerns interaction regimes that
are genuinely off-shell. Virtual, off-shell exchanges---van der Waals forces,
the Casimir effect, the Lamb shift---produce no DAG vertices in RA. They are
genuine physical effects, but they are effects of the potentia: they do not
constitute actualization events and do not contribute to decoherence.

This is not merely a restatement of existing results. The standard decoherence
account treats all environmental coupling as a potential source of decoherence,
modulated by coupling strength and spectral density. RA draws a qualitatively
different distinction: coupling that remains purely off-shell (because the
interaction energy does not reach the on-shell threshold) is decoherence-immune
in the perturbative radiative regime, not merely weakly decoherent. This holds
cleanly when on-shell boson emission is the dominant decoherence channel;
RAQM notes that near-field and many-body processes can produce environmental
distinguishability through other mechanisms, so the immunity is regime-specific
rather than categorical.

The potential design implication is significant. If a quantum computing
architecture can be engineered so that its dominant environmental coupling remains
below the on-shell threshold---so that qubit-environment interactions are
kinematically constrained to the virtual regime---RA predicts that this regime
would exhibit qualitatively suppressed decoherence, not merely quantitatively
reduced decoherence. This is a testable prediction that distinguishes RA from
standard decoherence theory.

The experimental platform most naturally positioned to probe this distinction is
the superconducting circuit, where the relevant on-shell bosons are microwave
photons and the threshold condition is set by the circuit's transition frequencies.
Operating in a regime where the environmental spectral density falls entirely below
the transition frequency creates exactly the off-shell condition RA identifies as
decoherence-immune. This is the physical interpretation of the ``bad cavity''
versus ``good cavity'' limit in circuit quantum electrodynamics.

\subsection{The Threshold Theorem Acquires a Physical Interpretation}

The fault-tolerance threshold theorem establishes that scalable quantum computation
is possible provided the physical error rate $p$ falls below a threshold $p_{\rm
th}$, which depends on the error correction code and architecture. The theorem is
currently stated in terms of error rates that are treated as phenomenological
parameters: there is no physical derivation of $p_{\rm th}$ from first principles.

If RA is correct, the physical error rate has a precise interpretation: it is the
rate of on-shell boson exchange between the logical qubit and its environment,
given by
\begin{equation}
  p \;\sim\; \frac{1}{\tau_{\mathrm{act}} \cdot f_{\mathrm{gate}}}
  \;=\; \frac{\Delta E_{\mathrm{env}}}{\hbar \cdot f_{\mathrm{gate}}},
\end{equation}
where $f_{\mathrm{gate}}$ is the gate operation frequency. The threshold
$p < p_{\rm th}$ then becomes a condition on the ratio
$\Delta E_{\mathrm{env}} / (\hbar \cdot f_{\mathrm{gate}})$: the environment must
not have sufficient energy to cause on-shell exchange faster than the gate cycle
time. This gives the threshold theorem a physical interpretation as a statement
about kinematic accessibility rather than an abstract error rate condition.

The further implication is that the threshold is not universal across
architectures in the way current theory suggests. For two architectures with the
same phenomenological error rate $p$ but different physical mechanisms---one
where errors arise from on-shell exchange and one where they arise from
control-pulse infidelity---RA predicts different scalability behaviour, because
only on-shell exchange constitutes genuine actualization. Control infidelity that
does not produce on-shell environmental exchange is, in principle, reversible and
correctable in a way that genuine actualization errors are not. In practice,
control pulse errors often couple to environmental modes that can go on-shell, so
this distinction is a conceptual one about error types rather than a claim about
the reversibility of all control errors in current hardware.

\subsection{The Kinematic Coherence Bound and the $1/N$ Law}
\label{sec:kcb}

The Kinematic Snap criterion yields a sharp, derivable bound on the maximum
achievable circuit depth for an $N$-qubit logical state, which we call the
\emph{Kinematic Coherence Bound} (KCB). The derivation proceeds directly from
Prediction~P2 of RAQM~\citep{Sandeman2026a} without additional assumptions.

\paragraph{Single-qubit baseline.}
For a single isolated qubit, the actualization timescale is
\begin{equation}
  \tau_{\mathrm{act}} \;=\; \frac{\hbar}{\Delta E_{\mathrm{env}}},
  \label{eq:tau1}
\end{equation}
where $\Delta E_{\mathrm{env}}$ is the characteristic energy of the environmental
modes available for on-shell exchange. The gate quality factor is
\begin{equation}
  \eta \;:=\; f_{\mathrm{gate}} \cdot \tau_{\mathrm{act}}
        \;=\; \frac{\hbar \, f_{\mathrm{gate}}}{\Delta E_{\mathrm{env}}},
  \label{eq:eta}
\end{equation}
where $f_{\mathrm{gate}}$ is the gate frequency. This is the RA physical
interpretation of the standard figure of merit $T_2 / t_{\mathrm{gate}}$: it is
the ratio of the gate period to the single-qubit actualization timescale.

\paragraph{$N$-qubit logical state.}
Consider an $N$-qubit register in an entangled logical state (for concreteness, a
GHZ-type state of the form $\frac{1}{\sqrt{2}}(\ket{00\cdots0} +
\ket{11\cdots1})$, or any state with genuinely $N$-body entanglement). The logical
information is encoded in correlations that are non-local across all $N$ qubits.

Under the Kinematic Snap criterion, each qubit constitutes an independent site at
which the kinematic threshold can be crossed: the actualization of any single qubit
destroys the logical entanglement of the entire register. The logical decoherence rate is therefore the sum of single-qubit actualization
rates across all $N$ sites:
\begin{equation}
  \Gamma_{\mathrm{logical}} \;=\; N \cdot \Gamma_1 \;=\; \frac{N \cdot \Delta E_{\mathrm{env}}}{\hbar},
  \label{eq:gamma_N}
\end{equation}
and the logical coherence time is
\begin{equation}
  \tau_{\mathrm{act}}^{(N)} \;=\; \frac{\hbar}{N \cdot \Delta E_{\mathrm{env}}}.
  \label{eq:tauN}
\end{equation}
This is the RA physical account of the well-known result from quantum information
theory that the coherence time of a GHZ state scales as $T_1/N$~\citep{Huelga1997}.
The Kinematic Snap mechanism explains \emph{why}: each qubit is a separate kinematic
threshold that can be independently crossed, and crossing any one of them destroys
the logical state.

\paragraph{The Kinematic Coherence Bound.}
The maximum circuit depth achievable before logical decoherence is
\begin{equation}
  \boxed{D_{\max} \;=\; f_{\mathrm{gate}} \cdot \tau_{\mathrm{act}}^{(N)}
                   \;=\; \frac{\eta}{N},}
  \label{eq:kcb}
\end{equation}
where $\eta$ is the single-qubit quality factor defined in~\eqref{eq:eta}. Circuit
depth scales strictly as $1/N$: one cannot increase the number of qubits and the
circuit depth simultaneously without improving $\eta$.

\paragraph{Maximum fault-tolerant array size.}
The fault-tolerance threshold condition requires the physical error rate per gate
to fall below a code-dependent threshold $p_{\mathrm{th}}$:
\begin{equation}
  p \;=\; \frac{\Gamma_{\mathrm{logical}}}{f_{\mathrm{gate}}}
    \;=\; \frac{N}{\eta} \;<\; p_{\mathrm{th}}.
  \label{eq:pth}
\end{equation}
This yields a \emph{maximum array size} for fault-tolerant operation at fixed
single-qubit quality:
\begin{equation}
  \boxed{N_{\max} \;=\; \eta \cdot p_{\mathrm{th}}.}
  \label{eq:nmax}
\end{equation}
This result is not anticipated by standard quantum error correction theory, which
holds that fault-tolerant quantum computation is scalable without bound once
$p < p_{\mathrm{th}}$ for individual qubits. The RA account adds a constraint:
the logical error rate grows with $N$ because each added qubit is an additional
kinematic site at which actualization can occur. Unlimited scaling without
improving $\eta$ will eventually push the logical error rate above
$p_{\mathrm{th}}$ regardless of single-qubit fidelity.

Two caveats on the derivation. First, it assumes independent decoherence channels
per qubit. Correlated errors---crosstalk, collective decoherence, or correlated
on-shell emission events---can produce worse-than-independent scaling, so $N_{{\max}}$
is an upper bound on the array size under independent depolarizing noise, not a
tight limit for all error models. Second, $p_{{\mathrm{th}}}$ is code- and
error-model-dependent; the numerical estimates below use a surface code with
depolarizing noise.

\paragraph{Syndrome extraction and the KCB.}
A natural question concerns the role of quantum error correction itself: does a
stabilizer measurement on an ancilla qubit constitute a Kinematic Snap that
destroys the logical state? The answer is no, and the distinction is physically
precise. In a surface code, syndrome extraction works by entangling ancilla qubits
with the data qubits and then measuring the ancillae. The ancilla measurement
\emph{does} constitute an actualization event---the ancilla crosses the on-shell
threshold and its state is irreversibly recorded.

However, the logical information
is encoded \emph{non-locally} across many physical data qubits; it is not localized
at any single kinematic site. The ancilla actualization reveals error syndrome
information without collapsing the logical state, because the logical degree of
freedom has no support at the ancilla's location in the causal DAG.

The KCB applies to \emph{unintended} environmental actualization events---on-shell
exchanges between a data qubit and an uncontrolled environmental mode---not to the
designed syndrome extraction operations, which are intentional actualization events
that extract entropy without disturbing the logical encoding. The $1/N$ scaling
therefore reflects the rate of \emph{logical errors} from environmental coupling,
which grows with $N$ under independent depolarizing noise. Continuous QEC protocols
suppress this error rate below $p_{{\mathrm{th}}}$ but cannot eliminate the
underlying kinematic constraint: $N_{{\max}}$ is the array size at which the
environmental logical error rate, even before QEC correction, equals $p_{{\mathrm{th}}}$.

Beyond $N_{{\max}}$, QEC must run faster than the logical decoherence rate, which
requires improving $\eta$---consistent with the empirical observation that scaling
superconducting arrays requires simultaneous improvement of single-qubit quality.

\paragraph{Classical control and the KCB.}
A natural objection: can arbitrarily fast classical control systems, running syndrome
extraction at rates approaching the environmental decoherence rate, circumvent
$N_{\max}$? The answer is no, for a kinematic reason. The KCB bound on $N_{\max}$
reflects the rate at which environmental on-shell exchanges produce logical errors
in the \emph{data qubits}, not in the ancillae or classical control hardware.

Syndrome extraction and classical decoding can correct an error after it occurs,
but they cannot prevent an on-shell environmental actualization event from
writing an error vertex into the causal DAG of a data qubit. The logical error
rate scales as $N \cdot \Gamma_{\mathrm{env}}$ for $N$-qubit arrays regardless of
the speed of the classical feedback loop, because the classical controller cannot
observe and suppress an on-shell exchange \emph{before} it produces an
irreversible causal record. Improving the syndrome extraction rate improves the
probability of catching and correcting errors within a QEC cycle; it does not
reduce the rate at which errors are generated.

The only way to reduce the
error-generation rate is to reduce $\Delta E_{\mathrm{env}}$ (improve isolation)
or increase $f_{\mathrm{gate}}$ (improve $\eta$). This is precisely what
$N_{\max} = \eta \cdot p_{\mathrm{th}}$ encodes: $N_{\max}$ is set by the
single-qubit quality factor $\eta$, not by the speed of the classical controller.

Concretely: for a surface code with $p_{\mathrm{th}} \approx 10^{-2}$ under
independent depolarizing noise, and a contemporary superconducting qubit with
$\eta \approx 10^3$--$10^4$ (corresponding to $T_2 \sim 100\,\mu\mathrm{s}$ and
$t_{\mathrm{gate}} \sim 100\,\mathrm{ns}$), the RA bound gives
$N_{\max} \sim 10$--$100$ logical qubits at fixed $\eta$. This is illustrative
rather than a fit to experimental data; the observation that large superconducting
arrays require continual improvement of $\eta$ alongside qubit count is consistent
with but does not uniquely confirm the KCB.

\paragraph{The Entropy per Vertex ($\langle\Delta S_{\mathrm{vertex}}\rangle$).}
The Kinematic Coherence Bound provides a direct theoretical bridge to the thermodynamic entropy generated by actualization. In quantum error correction, the threshold $p_{\mathrm{th}}$ represents the maximum environmental error rate that the syndrome extraction cycle can actively correct---the exact point at which the entropy pumped out by the classical controller balances the entropy generated by environmental actualization. Because each uncorrectable logical error corresponds to an actualization event crossing the kinematic threshold, the mean relative entanglement entropy generated per actualization vertex, $\langle\Delta S_{\mathrm{vertex}}\rangle$, is strictly bounded by the fault-tolerance threshold in the scalable limit. In natural units, we identify the scaling relation $\langle\Delta S_{\mathrm{vertex}}\rangle \approx p_{\mathrm{th}}$. This identification ($p_{\mathrm{th}} \approx 10^{-2}$ for standard topological codes) provides the critical microscopic entropy value required to derive the macroscopic gravitational coupling constant in the companion paper RADM~\citep{Sandeman2026RADM}.

\paragraph{The optimal isolation window.}
A related prediction concerns the engineering of qubit isolation. The target
quantity is $\Delta E_{\mathrm{env}}$: the energy available to environmental modes
for on-shell exchange. Reducing $\Delta E_{\mathrm{env}}$ directly extends
$\tau_{\mathrm{act}}$ and improves $\eta$. However, the environment has multiple
channels for on-shell exchange---photonic (microwave), phononic, and through
two-level system (TLS) defects in the substrate---each with its own $\Delta
E_{\mathrm{env}, i}$. Improving isolation in one channel (e.g., by increasing
cavity Q to suppress photonic decay) does not suppress the others.

RA therefore predicts that there is an \emph{optimal isolation window} $Q^*$ for
each channel: the cavity Q at which photonic actualization timescale $\tau_{\rm
photo} = \hbar / \Delta E_{\rm photo}$ equals the actualization timescale of the
dominant non-photonic channel $\tau_{\rm other} = \hbar / \Delta E_{\rm other}$.
Beyond $Q^*$, further improvement of photonic isolation yields diminishing returns
because the dominant decoherence pathway shifts to the non-photonic channel. This
is measurable as a change in the $T_1$ scaling with cavity Q: below $Q^*$ the
scaling should follow the Purcell prediction; above $Q^*$ it should plateau.

\subsection{Implications for Quantum Algorithm Design}

At a higher level of abstraction, the Kinematic Snap suggests a distinction
between quantum algorithms that operate entirely within the off-shell regime---
where the computation consists of reversible potentia evolving under unitary
dynamics---and those that necessarily involve actualization events. All
conventional quantum algorithms involve measurement, which is a forced
actualization. But the computation preceding measurement may be structured to
avoid premature actualization.

This distinction is already implicit in the design of adiabatic quantum computing
and quantum annealing, where the system is evolved slowly to avoid non-adiabatic
transitions. RA offers the following reinterpretation of this design principle:
avoiding non-adiabatic transitions is avoiding premature on-shell exchange with
environmental modes, which would constitute unintended actualization events.
The adiabatic condition, in RA terms, is the condition that the evolution remains
in the off-shell regime throughout. This is an interpretational reframing rather
than a new prediction; it does not alter the standard adiabatic theorem or provide
new algorithmic guidance beyond what standard theory already implies.

% ─────────────────────────────────────────────────────────────────────────────
\section{Implications for Quantum Thermodynamics}
\label{sec:thermodynamics}

\subsection{Landauer's Principle Grounded in the Causal DAG}
\label{sec:landauer}

Landauer's principle~\citep{Landauer1961} states that the logical erasure of one
bit of information dissipates a minimum energy of $k_B T \ln 2$ as heat into the
environment. The principle is standardly derived from thermodynamic arguments
(the second law applied to the Szilard engine) or from quantum information theory
(the quantum data compression theorem). Both derivations establish the bound; 
neither identifies the physical mechanism by which it is enforced.

In RA, Landauer's principle has a direct mechanical derivation from the Kinematic
Snap~\citep{Sandeman2026RACI}.

\begin{theorem}[Landauer Cost of Actualization, RACI Theorem~4]
Each irreversible actualization event that eliminates $\Delta q$ bits of quantum
uncertainty---reducing the von Neumann entropy of the system by
$\Delta S_{\mathrm{vN}} = \Delta q \cdot \ln 2$---dissipates at minimum
\begin{equation}
  W_{\mathrm{act}} \;=\; \Delta q \cdot k_B T \cdot \ln 2.
  \label{eq:landauer}
\end{equation}
\end{theorem}

The proof in RACI proceeds directly from the Kinematic Snap criterion: an
actualization event writes a permanent vertex to the causal DAG, which requires
on-shell boson emission into the environment. The von Neumann entropy decrease
$\Delta S_{\mathrm{vN}}$ of the system must be transferred to the environment as
heat $\geq T \Delta S_{\mathrm{vN}} = k_B T \cdot \Delta q \cdot \ln 2$. The
bound~\eqref{eq:landauer} follows.

What RA adds to Landauer's principle beyond the standard derivation:

\textbf{A mechanism.} The energy is dissipated specifically at the moment the
mediating boson crosses the on-shell threshold. The dissipation is not a
continuous process but a discrete event with a precise physical timing:
$\tau_{\mathrm{act}} \sim \hbar / \Delta E_{\mathrm{env}}$ before the on-shell
threshold is crossed.

\textbf{A scope condition.} The Landauer bound applies only to \emph{genuine
actualization events}---irreversible on-shell exchanges that write DAG vertices.
Operations that remain in the virtual, off-shell regime do not write DAG vertices
and can, in principle, be reversed without paying the Landauer cost. This is
consistent with the reversibility of coherent quantum operations; RA identifies the
physical criterion that distinguishes reversible from irreversible operations.

A note on the Reeb-Wolf derivation: Reeb and Wolf~\citep{ReebWolf2014} established
a fully quantum Landauer bound showing the minimum dissipation holds for any
entropy-decreasing operation on a quantum system in contact with a thermal
reservoir, without reference to emission events. The RA account does not contradict
this result; rather, RA provides the physical mechanism \emph{by which} the Reeb-Wolf
bound is saturated. In RA, an entropy-decreasing operation that does not involve
on-shell exchange is not genuinely irreversible---it is a coherent operation that
can, in principle, be reversed by subsequent unitary evolution. The Reeb-Wolf bound
becomes binding precisely when the operation crosses the Kinematic Snap threshold
and produces an on-shell environmental record. The two accounts agree on the bound;
they differ on the physical story of why it holds.

\textbf{A timescale.} The Landauer bound tells you the minimum energy cost of
erasure. RA additionally tells you the minimum \emph{time} for erasure:
$\tau_{\mathrm{min}} \sim \hbar / \Delta E_{\mathrm{env}}$. These two together
constitute a Landauer bound in phase space rather than merely in energy: the
erasure must cost at least $\hbar$ of action in addition to $k_B T \ln 2$ of
energy. This stronger version of the bound is a direct consequence of the
Kinematic Snap mechanism.

\subsection{The Maxwell's Demon Paradox Dissolved}

The Maxwell's demon~\citep{Maxwell1872} paradox asks how a microscopic agent that
sorts fast and slow molecules can apparently violate the second law of
thermodynamics, and what physical mechanism restores it. The standard resolution
via Landauer's principle~\citep{Bennett1982} argues that the demon's memory
erasure, not the measurement, pays the thermodynamic cost. This resolution, while
correct, involves an information-theoretic postulate (the equivalence of physical
and logical entropy) that RA makes unnecessary.

In RA, the demon is a physical system embedded in the causal DAG. Every act of
``measurement'' by the demon is an actualization event: the demon's sensory
apparatus exchanges on-shell bosons with the gas molecules in order to determine
their velocities. Each such exchange writes a vertex to the DAG and pays the
Landauer cost $W_{\mathrm{act}} \geq k_B T \ln 2$ at the moment of measurement,
not at the moment of memory erasure.

The standard resolution delays the cost to erasure; RA moves the cost to
measurement. Both arrive at the same thermodynamic accounting, but through
different physical pictures. The RA picture is arguably more natural: the cost is
paid when the physical interaction occurs, not when a logical register is cleared.

More precisely: there is no gap in RA between the demon's information acquisition
and the physical cost of that acquisition. The demon's acquisition of knowledge
about a molecule's velocity \emph{just is} the on-shell boson exchange between the
demon's sensory apparatus and the molecule. The thermodynamic cost is paid at that
moment. There is no separate step of ``storing the information'' and ``later
erasing it.'' The storage and the cost are the same physical event.

This removes the apparently paradoxical temporal gap in the standard resolution.
In the standard account, the demon seemingly operates for free during the sorting
phase and pays later during erasure---a picture that requires careful argument to
prevent apparent violations of the second law during the sorting phase itself. In
RA, the cost is paid continuously and locally during the measurement interactions,
and the second law is enforced vertex by vertex as the DAG grows.

\paragraph{Weak measurements.}
The argument above applies to strong projective measurements, which constitute
unambiguous actualization events in RA. Weak measurements---continuous monitoring
that acquires partial information without fully collapsing the quantum state---are
a more subtle case. A weak measurement that does not produce an on-shell
environmental record is not, in RA, a full actualization event; it is a
pre-actualization interaction that keeps the system in the reversible potentia
regime. Whether and how the thermodynamic cost accumulates continuously during
weak measurement, and at what point the Kinematic Snap threshold is crossed, is
an open question that connects to the measurement back-action literature and
constitutes a natural target for the RA framework's extension to continuous
monitoring scenarios.

\subsection{The Fluctuation Theorems with a Precise Actualization Criterion}

The Jarzynski equality~\citep{Jarzynski1997} and Crooks fluctuation
theorem~\citep{Crooks1999} relate the distribution of work performed on a system
driven away from equilibrium to the free energy difference between initial and
final states:
\begin{align}
  \langle e^{-W/k_BT} \rangle &= e^{-\Delta F/k_BT} \quad \text{(Jarzynski)}, \\
  \frac{P(+W)}{P(-W)} &= e^{(W - \Delta F)/k_BT} \quad \text{(Crooks)}.
\end{align}
These theorems hold for systems driven arbitrarily far from equilibrium and have
been verified in numerous experiments. They are currently derived from the
assumption of microscopic reversibility (detailed balance) and the identification
of work with energy exchanges during the driving protocol.

RA provides a physical decomposition of the work distribution in these theorems.

An important caveat governs the following discussion: the standard work $W$ in the
Jarzynski and Crooks theorems is defined as the total energy exchanged with the
external driving agent during the protocol. If RA were to redefine $W$ to include
only on-shell exchange events---call it $W_{\rm RA}$---then $W_{\rm RA} \neq W$
in general, and there is no a priori reason the standard equalities hold for the
restricted quantity $W_{\rm RA}$ without rederivation. The analysis below therefore
proceeds differently: it uses the standard $W$ and treats RA as providing a
physical decomposition of the \emph{histories} that contribute to the work
distribution, without redefining $W$ itself.

\textbf{The work distribution $P(W)$ decomposes over actualization sequences.}
In RA, each trajectory contributing to $P(W)$ is a causal DAG history: a specific
sequence of on-shell exchange events. Virtual, off-shell exchanges occur within
each trajectory but do not constitute irreversible events. The standard $W$ for
each trajectory is the total energy exchange; RA adds the physical picture that
this energy is concentrated in the on-shell exchange events that write DAG vertices.

\textbf{The forward and reverse protocols differ in their DAG histories.}
In standard statistical mechanics, detailed balance is stated over phase space
trajectories. RA proposes that the more fundamental objects are causal DAG
histories. The Crooks theorem, in this picture, says the ratio of probabilities
of a DAG history and its time-reverse is $e^{(W-\Delta F)/k_BT}$---where $W$ is
the standard work, computed over the full trajectory including virtual exchanges.

\textbf{Protocols that remain close to off-shell have fewer DAG vertices} and
contribute disproportionately to the low-work tail of $P(W)$. This is the physical
picture of why quasi-static (nearly reversible) protocols dominate the Jarzynski
average: they generate fewer irreversible actualization events per unit time.

This decomposition suggests a research program rather than a completed result: the
Jarzynski equality, restated as a sum over actualization sequences with the correct
weighting, may be more tractable for discrete few-body quantum systems where the
DAG is a finite combinatorial object. Whether the equality holds for a
suitably defined $W_{\rm RA}$ (restricted to on-shell contributions) is an open
question whose answer would either strengthen or constrain the RA framework.

\subsection{The Arrow of Time as Structural, Not Statistical}

Standard statistical mechanics derives the thermodynamic arrow of time from the
second law, which is in turn grounded in the initial low-entropy state of the
universe. This derivation is probabilistic: it says that entropy-increasing
processes are overwhelmingly more probable than entropy-decreasing ones, given
typical initial conditions. The underlying microscopic dynamics remain
time-reversible.

RA provides a structural account of the arrow of time that does not rely on
initial conditions or probabilistic arguments. The causal DAG grows in one
direction only: new vertices are added; old vertices are never removed. This is
not a consequence of the second law; it is a consequence of the ontological
commitment that actualization is primitive and irreversible. The past is fixed
because it has been written into the DAG. The future is open because it has not.

This distinction has a precise thermodynamic consequence. In standard statistical
mechanics, a process that violates the second law is \emph{possible} (it is merely
improbable). In RA, a process that removes a causal DAG vertex is not merely
improbable---it is categorically excluded by the ontology. The arrow of time is not
a statistical asymmetry; it is a logical one.

For quantum thermodynamics, this matters because many results in the field depend
on the assumption of microscopic reversibility. RA does not violate this
assumption at the level of the quantum dynamics---the Schr\"{o}dinger equation is
still time-reversible---but it adds an irreversible layer: the writing of DAG
vertices. The interplay between these two levels (reversible quantum dynamics and
irreversible DAG inscription) is precisely the interplay between pre-actualization
and post-actualization regimes that RAQM characterises via Prediction~P3.

% ─────────────────────────────────────────────────────────────────────────────
\section{The Quantum Information Theory Connection}
\label{sec:qit}

\subsection{Relative Entropy as the Fundamental Actualization Criterion}

The quantum relative entropy $S(\rho \| \sigma) = \mathrm{Tr}[\rho(\log \rho -
\log \sigma)]$ plays a central role in both quantum information theory and the RA
actualization criterion. In RAGC~\citep{Sandeman2026RAGC}, the fundamental
criterion for an actualization event is stated in terms of relative entropy 
with
respect to the vacuum: an actualization occurs when an interaction produces a
strictly irreversible increase in $S(\rho \| \sigma_0)$, where $\sigma_0$ is the
Minkowski vacuum state.

This criterion is frame-independent by the following chain of results:
\begin{enumerate}[noitemsep]
  \item $S(U\rho U^\dagger \| U\sigma U^\dagger) = S(\rho \| \sigma)$ for any
        unitary $U$ --- proved using \texttt{log\_unitary\_conj}, which is a
        one-line corollary of \texttt{Matrix.cfc\_conj\_unitary} from the
        Lean-QuantumInfo library~\citep{Meiburg2025LQI}.
  \item The Minkowski vacuum is Poincar\'{e}-invariant:
        $U\sigma_0 U^\dagger = \sigma_0$ --- a QFT physics input.
  \item Therefore $S(U\rho U^\dagger \| \sigma_0) = S(\rho \| \sigma_0)$ ---
        the compiled theorem \texttt{frame\_independence} in
        \texttt{RA\_AQFT\_Proofs\_v10.lean}.
\end{enumerate}

This means the actualization criterion is a covariant scalar: it is the same in
all inertial and uniformly accelerated frames. The Rindler thermal bath, seen by
an accelerated observer, has $\Delta S_{\mathrm{rel}} = 0$---it is stationary
with respect to the vacuum and does not constitute a sequence of actualization
events. This is the compiled theorem \texttt{rindler\_stationarity}.

\subsection{The Data Processing Inequality and Causal Structure}

The data processing inequality---that relative entropy can only decrease under the
application of a completely positive trace-preserving (CPTP) map,
$S(\mathcal{E}(\rho) \| \mathcal{E}(\sigma)) \leq S(\rho \| \sigma)$---has a
natural RA interpretation. A CPTP map represents a physical process that may
include actualization events (on-shell exchanges). Each such event can only reduce
the distinguishability between states $\rho$ and $\sigma$ as seen from outside the
system. Information is lost from the perspective of the environment when a DAG
vertex is written; it is not created.

The RA formulation of the data processing inequality is therefore a statement about
causal DAG evolution: each actualization step can only reduce relative entropy, not
increase it. An increase in relative entropy $S(\rho \| \sigma_0)$---an
actualization event---necessarily involves a physical interaction, not merely a
CPTP map. This distinguishes the two directions of the inequality sharply:
decreasing relative entropy corresponds to physical processes that may or may not
involve actualization; increasing relative entropy (an actualization event) requires
specific physical conditions (on-shell exchange).

The Petz monotonicity theorem~\citep{Petz1986}, which establishes the data
processing inequality for quantum relative entropy, is formalised in
Lean-QuantumInfo~\citep{Meiburg2025LQI} as part of the Generalised Quantum
Stein's Lemma proof infrastructure. This is the mathematical substrate for the
RA AQFT proof program.

\subsection{Quantum Hypothesis Testing and the Actualization Criterion}

The Generalized Quantum Stein's Lemma~\citep{Meiburg2025LQI} establishes the
operational significance of quantum relative entropy in the context of quantum
resource theories: it determines the optimal error exponents for discriminating
between a quantum state and a set of free states. We propose the following RA
interpretation of this result: $S(\rho \| \sigma_0)$ quantifies how distinguishable
the current state $\rho$ is from the vacuum $\sigma_0$, which is precisely the
quantity whose increase constitutes an actualization event. The GQSL itself makes
no claim about RA; the connection is a proposed conceptual bridge, not a derived
result.

Under this interpretation, the Stein's Lemma suggests an operational picture of
the actualization threshold: an actualization event corresponds to the system
crossing from a regime where $\rho$ is statistically indistinguishable from
$\sigma_0$ (in the quantum hypothesis testing sense) to one where the
distinguishability is irreversibly established. The type-I and type-II error
exponents in the Stein's Lemma would then determine the precision with which
this threshold can be located experimentally.

This connection suggests a natural experimental program: quantum hypothesis testing
experiments that probe the vicinity of the actualization threshold could provide
direct tests of the RA criterion. Specifically, the sharp threshold predicted by
RAQM (Prediction~P3) would manifest as a sharp transition in the hypothesis
testing error exponents as the environmental coupling crosses the on-shell
threshold.

% ─────────────────────────────────────────────────────────────────────────────
\section{Empirical Predictions and Experimental Touchstones}
\label{sec:predictions}

The implications developed above are conditional on the correctness of the RA
framework. The following predictions, inherited from RAQM and RAGC, provide the
experimental tests.

\subsection{Predictions Relevant to Quantum Computing (P1--P3)}

\paragraph{P1: Finite actualization propagation velocity.}
The actualization front in a spatially extended many-body quantum system propagates
at a velocity bounded by the Lieb-Robinson speed~\citep{LiebRobinson1972}
$v_{\mathrm{LR}} \sim J\ell/\hbar$. A decoherence event at site $A$ should
produce a time delay $\Delta t \geq |r_A - r_B|/v_{\mathrm{LR}}$ before
off-diagonal coherences at site $B$ begin to collapse. Instantaneous collapse
models predict no such delay.

\paragraph{Connection to causal bandwidth.}
In the RA-CSG framework, $v_{\mathrm{LR}}$ is the lattice-scale expression
of the same causal bandwidth ceiling as $c = l_P/t_P$ in the relativistic
regime: both are the maximum rate at which actualization can propagate
through the causal graph. The Kinematic Coherence Bound threshold
$p_{\mathrm{th}} \approx 10^{-2}$ is not a property of specific error
correction codes but the per-link actualization probability at which
the causal graph transitions from locally 2-dimensional to locally
3-dimensional (the Poisson-CSG $\mu = 1$ threshold). Fault-tolerant
quantum computing requires $\mu > 1$ in the logical qubit's causal
neighbourhood: the threshold is the same Erd\H{o}s-R\'{e}nyi giant
component transition that governs galactic dimensional reduction in
RADM~\citep{Sandeman2026RADM} and the biological Causal Firewall
in RACI~\citep{Sandeman2026RACI}. The fault-tolerance threshold,
the speed of light, and the Causal Firewall are all the same
physics at different scales of the causal graph.

\paragraph{P2: Environment-dependent transition duration.}
The actualization timescale scales as $\tau_{\mathrm{act}} \sim \hbar / \Delta
E_{\mathrm{env}}$. Systematically varying $\Delta E_{\mathrm{env}}$ in a
controlled superconducting circuit experiment should produce a proportional
variation in transition duration. This is consistent with Minev et
al.~\citep{Minev2019} and Pazourek et al.~\citep{Pazourek2015}; the distinctive
prediction is the \emph{proportional} scaling.

\paragraph{P3: Sharp reversibility threshold.}
There is a sharp threshold between pre-actualization (reversible) and
post-actualization (irreversible) dynamics, determined by the environmental
entanglement entropy crossing a critical value. Identifying this threshold
experimentally constitutes a direct measurement of the actualization criterion.
Decoherence-only models predict no such sharp threshold.

\subsection{Predictions Relevant to Quantum Thermodynamics}

\paragraph{P4 (from P2): Landauer timescale.}
In addition to the Landauer energy bound $W_{\mathrm{act}} \geq k_B T \ln 2$, RA
predicts a Landauer \emph{timescale}: erasure cannot be completed faster than
$\tau_{\mathrm{min}} \sim \hbar / \Delta E_{\mathrm{env}}$. This sets a joint
energy-time bound on erasure that is stronger than the standard Landauer bound and
is, in principle, measurable in fast qubit experiments.

\paragraph{P5: Off-shell regime is thermodynamically inert.}
Interactions that remain kinematically constrained to the virtual, off-shell regime
do not dissipate energy in the thermodynamic sense. This is testable in carefully
engineered systems where the dominant qubit-environment coupling can be tuned below
the on-shell threshold, predicting a qualitative suppression of thermodynamic
dissipation rather than a quantitative reduction.

% ─────────────────────────────────────────────────────────────────────────────
\section{Discussion}
\label{sec:discussion}

\subsection{What RA Provides That Existing Frameworks Do Not}

The implications developed in Sections~\ref{sec:computing}--\ref{sec:qit} share a
common structure: RA provides a \emph{mechanism} where existing frameworks provide
only \emph{statistics}. Standard decoherence theory tells you the rate at which
coherence is lost; RA tells you the physical process by which each individual
coherence loss occurs, and gives it a precise timing. Standard quantum
thermodynamics invokes entropy and irreversibility; RA grounds these in the
discrete, irreversible growth of a causal DAG.

This shift from statistical to mechanical description is not merely aesthetic.
It has three practical consequences:

First, it gives engineering targets that are physical rather than phenomenological.
The target is $\Delta E_{\mathrm{env}}$---a quantity with a direct physical
interpretation---rather than $T_2$---a phenomenological timescale whose physical
determinants must be inferred separately.

Second, it identifies qualitative distinctions (off-shell versus on-shell) that the
statistical description treats as quantitative ones (strong coupling versus weak
coupling). The difference between a system that is decoherence-immune because it
is off-shell and a system that is weakly decohering because it is weakly coupled is
qualitative in RA and quantitative in standard theory. That qualitative distinction
may be experimentally accessible.

Third, it connects the two fields to each other and to the broader RA framework.
Quantum computing and quantum thermodynamics share the relative entropy as their
central quantity; RA gives that quantity a physical interpretation as the
actualization criterion; and the frame-independence of that criterion is a compiled
Lean theorem. The mathematical infrastructure is shared.

\subsection{Scope and Limitations}

The implications in this paper are derived under the assumption that the RA
actualization criterion is correct. That assumption is currently under peer review
in the companion paper RAQM~\citep{Sandeman2026a}. Several implications---
particularly the qualitative decoherence-immunity of the off-shell regime
(Section~\ref{sec:computing}) and the Landauer timescale bound (P4)---go beyond
what existing experimental data require and constitute genuine predictions that may
falsify the framework if tested and found wanting. This is the appropriate
scientific posture; the implications are stated as conditionals, not assertions.

The Lean-verified results cited in Section~\ref{sec:qit} establish the
mathematical consistency of the AQFT proof program in finite dimensions, with three
axioms: the CFC unitary-conjugation lemma (sourced from Lean-QuantumInfo), the
Poincar\'{e}-invariance of the vacuum (a QFT physics input), and the data
processing inequality (proved in Lean-QuantumInfo). The extension to the full
infinite-dimensional type~III$_1$ algebra setting remains the primary open
mathematical target.

% ─────────────────────────────────────────────────────────────────────────────
\section{Conclusion}

Relational Actualism proposes that the precise physical mechanism of quantum
irreversibility is the crossing of the on-shell dispersion relation by a mediating
boson. This paper has developed the implications of this proposal for quantum
computing and quantum thermodynamics. If correct, RA gives both fields what they
currently lack: a precise physical account of when and why quantum processes become
irreversible.

For quantum computing: the engineering target is the environmental energy
$\Delta E_{\mathrm{env}}$; the off-shell regime is categorically
decoherence-immune; the threshold theorem acquires a kinematic interpretation; and
the distinction between reversible and irreversible errors is physically grounded.

For quantum thermodynamics: Landauer's principle has a mechanical derivation with
an additional timescale bound; the Maxwell's demon pays its thermodynamic cost at
measurement, not at erasure; the fluctuation theorems are restated as sums over
actualization sequences; and the arrow of time is structural rather than
statistical.

The shared mathematical substrate---quantum relative entropy, CPTP maps, the data
processing inequality---connects both fields to the RA AQFT proof program and to
the Lean-QuantumInfo library~\citep{Meiburg2025LQI} in which key results are
formally verified. The frame-independence of the actualization criterion is a
compiled theorem. The remaining mathematical frontier is the extension from finite
dimensions to the type~III$_1$ von Neumann algebra setting.

The framework makes specific, discriminating predictions (Section~\ref{sec:predictions})
that are testable with existing or near-future experimental platforms. Those tests
are the appropriate next step.

% ─────────────────────────────────────────────────────────────────────────────
\begin{thebibliography}{99}

\bibitem[Bennett(1982)]{Bennett1982}
C.~H.~Bennett.
\newblock The thermodynamics of computation---a review.
\newblock \textit{International Journal of Theoretical Physics}, 21:905--940,
  1982.

\bibitem[Crooks(1999)]{Crooks1999}
G.~E.~Crooks.
\newblock Entropy production fluctuation theorem and the nonequilibrium work
  relation for free energy differences.
\newblock \textit{Physical Review E}, 60(3):2721--2726, 1999.


\bibitem[Huelga et~al.(1997)]{Huelga1997}
S.~F.~Huelga, C.~Macchiavello, T.~Pellizzari, A.~K.~Ekert, M.~B.~Plenio, and
  J.~I.~Cirac.
\newblock Improvement of frequency standards with quantum entanglement.
\newblock \textit{Physical Review Letters}, 79(20):3865--3868, 1997.

\bibitem[Jarzynski(1997)]{Jarzynski1997}
C.~Jarzynski.
\newblock Nonequilibrium equality for free energy differences.
\newblock \textit{Physical Review Letters}, 78(14):2690--2693, 1997.

\bibitem[Landauer(1961)]{Landauer1961}
R.~Landauer.
\newblock Irreversibility and heat generation in the computing process.
\newblock \textit{IBM Journal of Research and Development}, 5(3):183--191, 1961.

\bibitem[Lieb and Robinson(1972)]{LiebRobinson1972}
E.~H.~Lieb and D.~W.~Robinson.
\newblock The finite group velocity of quantum spin systems.
\newblock \textit{Communications in Mathematical Physics}, 28(3):251--257, 1972.

\bibitem[Maxwell(1872)]{Maxwell1872}
J.~C.~Maxwell.
\newblock \textit{Theory of Heat}.
\newblock Longmans, Green and Co., London, 4th edition, 1872.

\bibitem[Meiburg et~al.(2025)]{Meiburg2025LQI}
A.~Meiburg, L.~A.~Lessa, and R.~R.~Soldati.
\newblock A formalization of the generalized quantum {Stein's} lemma in {Lean}.
\newblock \textit{arXiv}:2510.08672, 2025.
\newblock Lean-QuantumInfo library:
  \url{https://github.com/Timeroot/Lean-QuantumInfo}.

\bibitem[Minev et~al.(2019)]{Minev2019}
Z.~K.~Minev, S.~O.~Mundhada, S.~Shankar, P.~Reinhold, R.~Guti\'{e}rrez-J\'{a}uregui,
  R.~J.~Schoelkopf, M.~Mirrahimi, H.~J.~Carmichael, and M.~H.~Devoret.
\newblock To catch and reverse a quantum jump mid-flight.
\newblock \textit{Nature}, 570:200--204, 2019.

\bibitem[Pazourek et~al.(2015)]{Pazourek2015}
R.~Pazourek, S.~Nagele, and J.~Burgd\"{o}rfer.
\newblock Attosecond chronoscopy of photoemission.
\newblock \textit{Reviews of Modern Physics}, 87(3):765--802, 2015.

\bibitem[Petz(1986)]{Petz1986}
D.~Petz.
\newblock Sufficient subalgebras and the relative entropy of states of a von
  {Neumann} algebra.
\newblock \textit{Communications in Mathematical Physics}, 105(1):123--131,
  1986.


\bibitem[Reeb and Wolf(2014)]{ReebWolf2014}
D.~Reeb and M.~M.~Wolf.
\newblock An improved {Landauer} principle with finite-size corrections.
\newblock \textit{New Journal of Physics}, 16(10):103011, 2014.

\bibitem[Sandeman(2026e)]{Sandeman2026RAHC}
J.~F.~Sandeman.
\newblock Algorithmic Causal Dynamics: A Unified Framework for Emergence and
  Complexity in Relational Actualism ({RAHC}).
\newblock Preprint, 2026.
\newblock \href{https://doi.org/10.5281/zenodo.19198171}{doi.org/10.5281/zenodo.19198171}.

\bibitem[Sandeman(2026a)]{Sandeman2026a}
J.~F.~Sandeman.
\newblock Relational Actualism and Quantum Mechanics ({RAQM}).
\newblock \textit{Foundations of Physics} (Under review), 2026.
\newblock \href{https://doi.org/10.5281/zenodo.19174380}{doi.org/10.5281/zenodo.19174380}.

\bibitem[Sandeman(2026b)]{Sandeman2026RAGC}
J.~F.~Sandeman.
\newblock Relational Actualism: Gravity and Cosmology ({RAGC}).
\newblock Preprint, 2026.
\newblock \href{https://doi.org/10.5281/zenodo.19197999}{doi.org/10.5281/zenodo.19197999}.

\bibitem[Sandeman(2026c)]{Sandeman2026RACI}
J.~F.~Sandeman.
\newblock Relational Actualism: Complexity and Intelligence ({RACI}).
\newblock Preprint, 2026.
\newblock \href{https://doi.org/10.5281/zenodo.19198224}{doi.org/10.5281/zenodo.19198224}.

\bibitem[Bland-Hawthorn \& Gerhard(2016)]{BlandHawthorn2016}
J.~Bland-Hawthorn and O.~Gerhard.
\newblock The Galaxy in context: Structural, kinematic, and integrated
  properties.
\newblock \emph{Annual Review of Astronomy and Astrophysics}, 54:529--596, 2016.
\newblock \href{https://doi.org/10.1146/annurev-astro-081915-023441}{doi:10.1146/annurev-astro-081915-023441}.

\bibitem[Sandeman(2026f)]{Sandeman2026RADM}
J.~F.~Sandeman.
\newblock Relational Actualism: Dark Matter as Actualization-Density Topology ({RADM}).
\newblock Preprint, 2026.
\newblock \href{https://doi.org/10.5281/zenodo.19198513}{doi.org/10.5281/zenodo.19198513}.

\end{thebibliography}

\section*{Declarations}
\textbf{Funding:} Not applicable.\quad
\textbf{Competing interests:} The author declares no competing interests.

\textbf{Preprint availability:} This manuscript is available as a preprint at
\href{https://doi.org/10.5281/zenodo.19198285}{doi.org/10.5281/zenodo.19198285}.

\textbf{Code availability:} The Lean~4 formal proof file (\texttt{RA\_AQFT\_Proofs\_v10.lean}) and Python computational scripts are publicly available at \href{https://github.com/jsandeman/Relational-Actualism}{github.com/jsandeman/Relational-Actualism}.

\textbf{Data availability:} This paper presents theoretical work. No datasets
were generated or analysed.

\end{document}